\documentclass[11pt,a4paper]{article}
% ---------------------------------------------------------------------------
%  Summa -- the style of the showcase and of the manual: the common style of the packages, n4sd4q.sty (a copy of
%  the one of the hub, with the palette n4sd4q-palette.tex: found through TEXINPUTS, set by build_pdf.sh and
%  build_manual.sh), and what is particular to Summa. The examples are computed by Development/run_examples.wl:
%  examples/NAME.wl is the code, NAME.tex the output as a formula, NAME.txt the output as text, NAME.pdf the output
%  as a picture (a diagram, a contraction tree, a plot, the summary box of an EinsumFunction).
% ---------------------------------------------------------------------------
\usepackage[utf8]{inputenc}
\usepackage{n4sd4q}
\familysetup{package=Summa, logo=../Logo/Summa_logo_white.pdf}

\familykeywords{Einsum,EinsumPath,EinsumFunction,EinsumDiagram,EinsumTree,EinsumClearCache,SummaUpdateCheck}
\lstset{literate={α}{{$\alpha$}}1 {β}{{$\beta$}}1 {γ}{{$\gamma$}}1 {δ}{{$\delta$}}1 {ε}{{$\varepsilon$}}1 {ϵ}{{$\epsilon$}}1 {η}{{$\eta$}}1
  {μ}{{$\mu$}}1 {ν}{{$\nu$}}1 {σ}{{$\sigma$}}1 {τ}{{$\tau$}}1 {ψ}{{$\psi$}}1 {φ}{{$\varphi$}}1 {ϕ}{{$\phi$}}1 {χ}{{$\chi$}}1
  {ω}{{$\omega$}}1 {→}{{$\to$}}1 {…}{{$\ldots$}}1 {·}{{$\cdot$}}1 {∑}{{$\sum$}}1 {×}{{$\times$}}1}
% the same letters in the text (\texttt{α[k]}, as the code writes them)
\DeclareUnicodeCharacter{03B1}{\ensuremath{\alpha}}\DeclareUnicodeCharacter{03B2}{\ensuremath{\beta}}
\DeclareUnicodeCharacter{03C3}{\ensuremath{\sigma}}\DeclareUnicodeCharacter{2192}{\ensuremath{\to}}

\newcommand{\exmissing}[1]{\textcolor{red}{[missing example: \detokenize{#1}]}}
\makeatletter
\newcommand{\excodeafter}{\tcb@parfillskip@check\par\ifvmode\nobreak\@nobreaktrue\nointerlineskip\addvspace{\glueexpr2pt-\parskip}\fi}
\makeatother
\newcommand{\excode}[1]{\ifhmode{\widowpenalty=10000\par}\nobreak\fi\IfFileExists{examples/#1.wl}{\tcbinputlisting{listing only,
  colback=codebg, colframe=codeframe, boxrule=0.5pt, arc=2pt, left=6pt, right=6pt, top=3pt, bottom=3pt, before skip=4pt,
  after=\excodeafter, listing file={examples/#1.wl}, listing options={language=WL}}}{\exmissing{#1}}}
\newcommand{\exformula}[1]{\begin{exoutput}\raggedright\IfFileExists{examples/#1.tex}{$\displaystyle\input{examples/#1.tex}$}{\exmissing{#1}}\end{exoutput}}
\newcommand{\extext}[1]{\IfFileExists{examples/#1.txt}{\tcbinputlisting{listing only, colback=white, colframe=codeframe,
  boxrule=0.5pt, arc=2pt, left=24pt, right=6pt, top=2pt, bottom=2pt, before skip=2pt, after skip=8pt,
  listing file={examples/#1.txt}, listing options={style=plainoutput}, overlay unbroken and first={\node[anchor=north west,
  font=\scriptsize\sffamily\color{comment}] at ([xshift=4pt, yshift=-5pt]frame.north west) {out};}}}{\exmissing{#1}}}
% a picture, centred in the box of the output, on white: at the size of its PDF (the front end exports it at 72/100 of
% its size on the screen), or as wide as the box if it is wider (\exprepare{width}{NAME} sets the width instead). The code
% and the picture stay on the same page
\newsavebox{\expicturebox}
\newcommand{\exprepare}[2]{\ifx\relax#1\relax\sbox{\expicturebox}{\includegraphics{examples/#2.pdf}}%
  \ifdim\wd\expicturebox>\dimexpr\linewidth-38pt\relax
  \sbox{\expicturebox}{\resizebox{\dimexpr\linewidth-38pt\relax}{!}{\includegraphics{examples/#2.pdf}}}\fi
  \else\sbox{\expicturebox}{\includegraphics[width=#1]{examples/#2.pdf}}\fi}
\newcommand{\expicture}{\begin{exoutput}\centering\usebox{\expicturebox}\end{exoutput}}
% room for a picture and its code: a new page when the rest of this one is too short. Unlike \needspace, no negative
% penalty is left when the room is there, so that the page is not broken here anyway; \keepfor[lines]{NAME} before the
% text that leads to a picture (headings included, about so many lines with the code) keeps the text with it
\makeatletter
\newcommand{\exkeepspace}[1]{\par\penalty\@M\begingroup\dimen@=\dimexpr#1\relax\dimen@ii=\pagegoal\advance\dimen@ii-\pagetotal
  \ifdim\dimen@>\dimen@ii\ifdim\dimen@ii>\z@\vfil\penalty-\@M\fi\fi\endgroup}
\makeatother
\newcommand{\expicturekeep}[2]{\IfFileExists{examples/#2.pdf}{\exprepare{#1}{#2}%
  \exkeepspace{\ht\expicturebox+\dp\expicturebox+8\baselineskip}}{}}
\newcommand{\keeplines}[1]{\exkeepspace{#1\baselineskip}}   % room for a table and the sentence before it
\newcommand{\keepfor}[2][8]{\IfFileExists{examples/#2.pdf}{\exprepare{}{#2}%
  \exkeepspace{\ht\expicturebox+\dp\expicturebox+#1\baselineskip+3\baselineskip}}{}}   % [lines of text and code before it]
\newcommand{\example}[1]{\expicturekeep{}{#1}\excode{#1}\IfFileExists{examples/#1.tex}{\exformula{#1}}{\IfFileExists{examples/#1.txt}{\extext{#1}}%
  {\IfFileExists{examples/#1.pdf}{\expicture}{}}}}
% an example whose picture is drawn at a given width
\newcommand{\examplewidth}[2]{\expicturekeep{#1}{#2}\excode{#2}\IfFileExists{examples/#2.pdf}{\expicture}{\exmissing{#2}}}

\familysetup{document=Reference manual}
\newcommand{\fnentry}[2]{\par\needspace{4\baselineskip}\noindent\fn{#1}\par\nobreak
  \noindent\begin{minipage}[t]{\linewidth}\leftskip=1.2em #2\end{minipage}\par}

\begin{document}
\familytitle{Reference manual}{Every function, option and convention, with an example computed by the package}{version 1.0.0 \textperiodcentered{} November 2026}

{\small\tableofcontents}

\section{Installation and loading}
Summa is a paclet. Install the paclet of the latest \href{https://github.com/n4sd4q/Summa/releases}{release} once, then load
it in every session; the guided tour comes with the paclet:
\begin{tcblisting}{listing only, colback=codebg, colframe=codeframe, boxrule=0.5pt, arc=2pt, left=6pt, right=6pt, top=3pt, bottom=3pt,
  listing options={language=WL, basicstyle=\ttfamily\fontsize{7.5pt}{9.6pt}\selectfont}}
PacletInstall["https://github.com/n4sd4q/Summa/releases/download/v1.0.0/Summa-1.0.0.paclet"]
Needs["Summa`"]
NotebookOpen[PacletObject["Summa"]["AssetLocation", "Examples"]]
\end{tcblisting}
To use a clone of the repository without installing, load its folder before \texttt{Needs}:
\begin{code}
PacletDirectoryLoad["/path/to/Summa/Summa"]
\end{code}
Summa has no dependencies. It needs Mathematica 12.1 or later and is tested with 15.0. In the examples of this manual it is
loaded by
\example{load}
\fnentry{\$SummaVersion}{The version of the package, a string.}
\example{version}
Loaded in a notebook, the package prints a banner: its logo, the attribution ``Summa, by Nazzareno Africani'', the version,
the license, how to cite it, and links to the repository and the website. A script prints nothing and never checks for
updates.
\fnentry{SummaUpdateCheck, \$SummaUpdateCheck}{At most once a day, in the background, loading the package in a notebook asks
GitHub (\texttt{api.github.com}) for the latest release and, if it is newer than \texttt{\$SummaVersion}, adds to the banner
the \texttt{PacletInstall} line of that release; any failure is silent, and a session without access to the Internet never
asks. \texttt{\$SummaUpdateCheck = False} before \texttt{Needs["Summa\textasciigrave"]} (in \texttt{init.m} for good) turns the check off;
the setting is kept when the package is loaded again. \fn{SummaUpdateCheck[]} asks by hand and waits for the answer:
it gives \texttt{<|"InstalledVersion" -> \dots, "LatestVersion" -> \dots, "InstallLine" -> \dots|>}, and prints the
\texttt{PacletInstall} line when the latest release is newer; when the release cannot be read it gives \texttt{\$Failed}
with the message \texttt{SummaUpdateCheck::latest}. This is the only access of the package to the network.}
The test suite runs with \texttt{wolfram -script Tests/RunTests.wls}.

\section{Conventions}
\begin{itemize}
\item \textbf{Subscripts.} In a string specification each index is one character: letters, digits or any Unicode letter,
  case-sensitive (\texttt{i} and \texttt{I} are different indices). Subscripts are separated by commas; spaces, tabs and new
  lines are ignored. The arrow is \texttt{->}, \texttt{→} (\texttt{\textbackslash[RightArrow]}), \texttt{\textbackslash[LongRightArrow]}
  or \texttt{\textbackslash[Rule]}. In the list syntax a subscript is a list, and an index is any expression.
\item \textbf{Summation.} An index that appears in the operands and not in the output is summed over; an index of the output
  is kept. An index can appear any number of times, within one operand (diagonals and traces) and across operands (batch
  indices and hyperedges). An index of the output cannot repeat, as in numpy.
\item \textbf{Implicit output.} Without the arrow, the output has the indices that appear exactly once, sorted: by character
  code in a string specification (upper case before lower case), in canonical order (\texttt{Sort}) in the list syntax;
  the indices of the ellipsis come first.
\item \textbf{Ellipsis.} \texttt{...} (or \texttt{…}) stands for the leading axes that a subscript does not name, at most once
  per subscript. The ellipses of the operands are aligned from the right and broadcast as in numpy; when the operands have
  such axes, an explicit output must contain \texttt{...}.
\item \textbf{Dimensions.} An index has one dimension in all its places; an axis of size 1 is broadcast against a larger
  one. The places of an index repeated within one operand must have the same dimension.
\item \textbf{Operands.} A tensor is a rectangular list (packed or not) or a \texttt{SparseArray}; any other argument,
  \texttt{Sqrt[2]} or \texttt{a + b} included, is a scalar, with the empty subscript. Entries can be numbers of any kind or
  symbolic expressions. A result with no index is a scalar.
\item \textbf{Paths.} A contraction path is a list of pairs in the linear format of numpy and opt\_einsum, 1-based: each pair
  gives the positions of two tensors in the current list, which are removed; their result is appended at the end.
\item \textbf{Costs.} The cost of a pairwise contraction is the number of terms of its sum, the product of the dimensions of
  all its indices; the cost of a path (\opt{Flops}) is the sum over its steps, plus the size of each operand with a
  diagonal or an index summed alone. The naive cost (\opt{NaiveFlops}) is the product of the dimensions of all the indices
  times the number of operands minus one (at least one).
\end{itemize}
Spaces, Unicode letters and the arrow $\to$; and the case of the indices:
\example{labels}
\example{casesensitive}
The implicit output, in a string and in the list syntax (\opt{Spec} of \fn{EinsumPath} writes it out):
\example{implicitorder}
The ellipsis comes first in the implicit output; it can also stand for trailing axes:
\example{ellipsisimplicit}
Broadcasting of an axis of size 1, and the dimensions of a diagonal, which are not broadcast:
\example{broadcastrule}
\example{diagonalsame}
Scalars, and a sum over all the entries:
\example{scalar}

\section{Einsum}
\fnentry{Einsum}{\fn{Einsum["spec", t1, t2, \dots]} evaluates the einsum of the string \texttt{spec}, as
\texttt{"ij,jk->ik"}. \fn{Einsum[\{sub1, sub2, \dots\} -> out, t1, t2, \dots]} takes the subscripts as lists of indices
(or as strings), \fn{Einsum[\{sub1, sub2, \dots\}, \dots]} has the implicit output and \fn{Einsum["in" -> "out", \dots]} is the
string form written as a rule. The order of the contractions is chosen, each pairwise contraction runs on a kernel (BLAS
\texttt{Dot}, a batched or a broadcasting product), and the plan is cached.}
\example{forms}
\keeplines{12}Options:

{\small\begin{tabularx}{\linewidth}{llX}\toprule option & default & values \\ \midrule
\opt{Path} & \texttt{Automatic} & \texttt{Automatic}: \opt{Optimal} up to 8 tensors, \opt{Greedy} beyond; \opt{Optimal}: the
 cheapest path, by dynamic programming over the subsets of the tensors (up to 14 tensors); \opt{Greedy}: at each step the
 pair that least increases the size (opt\_einsum); \opt{Sequential}: from left to right; \texttt{\{\{i, j\}, \dots\}}: a
 complete path of your own (linear format, 1-based) \\
\opt{BatchMethod} & \texttt{Automatic} & the kernel of a batched contraction: \texttt{Automatic} (a cost model), \opt{Dot} (a
 loop of \texttt{Dot}) or \opt{Vectorized} (products and \texttt{Total} on packed arrays) \\
\opt{Cache} & \texttt{True} & \texttt{False}: the plan is made again and not kept \\ \bottomrule\end{tabularx}}

\smallskip
The result does not depend on the options, only the time to get it:
\example{optpath}
\example{optbatch}
\example{optcache}
An unknown option gives a message and is ignored:
\example{optunknown}

\section{EinsumPath}
\fnentry{EinsumPath}{\fn{EinsumPath[spec, t1, t2, \dots]} or \fn{EinsumPath[spec, \{dims1, dims2, \dots\}]} describes the
contraction path of an einsum and its cost, from the shapes alone: an association with}

{\small\begin{tabularx}{\linewidth}{lX}\toprule key & \\ \midrule
\opt{Spec} & the specification, written out (with the implicit output) \\
\opt{Path} & the path, in the linear format of numpy, 1-based \\
\opt{Method} & \opt{Optimal}, \opt{Greedy}, \opt{Sequential}, \opt{User}, or \opt{Trivial} (one or two tensors) \\
\opt{Contractions} & the specification of each pairwise contraction \\
\opt{StepFlops} & the cost of each pairwise contraction \\
\opt{Flops} & the cost of the path \\
\opt{NaiveFlops} & the cost of the sum written out over all the indices at once \\
\opt{Speedup} & \opt{NaiveFlops}/\opt{Flops} \\
\opt{LargestIntermediate} & the size of the largest tensor built on the way \\ \bottomrule\end{tabularx}}

\smallskip
Option \opt{Path}, with the values of \fn{Einsum}. The product of three matrices: $A$ times $B$ first, which is cheaper
than $B$ times $C$ first:
\example{path}
\example{pathuser}
The tensors or their dimensions give the same description:
\example{pathtensors}
With one or two tensors there is nothing to choose; an operand with a diagonal or an index summed alone costs its size,
and the speed-up can then be below 1:
\example{pathtrivial}

\section{EinsumFunction}
\keepfor[10]{efsummary}\fnentry{EinsumFunction}{\fn{EinsumFunction[spec, \{dims1, dims2, \dots\}]} makes the plan of an einsum once, for tensors of
the given dimensions (\texttt{\{\}} for a scalar), and gives a function: \fn{ef[t1, t2, \dots]} runs the plan. Tensors of other
dimensions give \texttt{Einsum::mismatch} and \texttt{\$Failed}. Options \opt{Path} and \opt{BatchMethod}, as \fn{Einsum}. In
a notebook it shows a summary box with the icon of the network, the specification and the cost; opened, it shows the
shapes, the naive cost, the speed-up, the largest intermediate and the path.}
\example{efsummary}
\example{efapply}
\example{efscalar}
\keeplines{14}Its properties:

{\small\begin{tabularx}{\linewidth}{lX}\toprule property & \\ \midrule
\opt{Spec}, \opt{Path}, \opt{Method}, \opt{Flops}, \dots & \opt{NaiveFlops}, \opt{Speedup}, \opt{LargestIntermediate},
 \opt{Contractions}: as in \fn{EinsumPath} \\
\opt{Shapes} & the dimensions of the operands \\
\opt{Diagram} & \fn{EinsumDiagram} of the einsum, with the dimensions and \texttt{"ShowPath" -> True} \\
\opt{Tree} & \fn{EinsumTree} of the einsum, with the dimensions \\
\opt{Plan} & the plan itself, an association, for inspection (its form may change between versions) \\
\opt{Properties} & the list of the properties \\ \bottomrule\end{tabularx}}
\example{efproperties}
\example{efvalues}
\example{efplan}
\example{efdiagram}
\example{efmismatch}

\section{EinsumDiagram}
\fnentry{EinsumDiagram}{\fn{EinsumDiagram[spec]} draws the tensor network of an einsum; \fn{EinsumDiagram[spec, t1, t2, \dots]}
or \fn{EinsumDiagram[spec, \{dims1, dims2, \dots\}]} also writes the dimensions. Every index is a COPY tensor with one leg for
each time it appears, plus one if it is in the output, and is drawn by the number of its legs:}

{\small\begin{tabularx}{\linewidth}{lX}\toprule index & drawn as \\ \midrule
two legs, in two tensors & a bond, with the name of the index in the middle \\
two legs, in one tensor and the output & an open leg, in blue \\
two legs, in one tensor & a loop: a trace \\
one leg (summed in one tensor alone) & a dashed leg ending in a small circle: a contraction with the vector of ones \\
three legs or more & a hyperedge: a purple dot joined to the tensors (and an open leg if it is in the output) \\ \bottomrule
\end{tabularx}}

\smallskip
Bonds get thicker as their dimension grows (on a logarithmic scale); tensors with the same name have the same colour;
in a notebook, a tooltip on each element tells its indices, dimensions and role. The automatic layout is classical
multidimensional scaling of the distances in the graph, refined by stress majorization, with the dangling legs in the free
angular gaps. Options:

{\small\begin{tabularx}{\linewidth}{llX}\toprule option & default & \\ \midrule
\opt{TensorNames} & \texttt{Automatic} & the names of the symbols passed as tensors, otherwise A, B, C, \dots; or a list \\
\opt{ShowDimensions} & \texttt{Automatic} & the dimensions next to the indices, when they are known; \texttt{False} \\
\opt{ShowPath} & \texttt{False} & \texttt{True}: each bond marked with the step at which it is contracted \\
\opt{Path} & \texttt{Automatic} & the path drawn by \opt{ShowPath} and \opt{Steps}, as in \fn{Einsum} \\
\opt{Layout} & \texttt{Automatic} & \opt{Circle}, \opt{Line}, or the coordinates of the tensors, \texttt{\{\{x1, y1\}, \dots\}} \\
\opt{Highlight} & \texttt{\{\}} & the positions of the tensors drawn with a thick border \\
\opt{Steps} & \texttt{False} & \texttt{True}: a list of pictures, one per step and one of the result \\
\opt{DefaultDimension} & 8 & the dimension of every index for the path, when no dimensions are given \\
\texttt{ImageSize}, \texttt{Background} & \texttt{Automatic}, \texttt{White} & as in \texttt{Graphics} \\ \bottomrule\end{tabularx}}

\smallskip
\keepfor[4]{diagramplain}A hyperedge kept in the output:
\example{diagramplain}
\keepfor[4]{diagramcircle}The layouts \opt{Circle}, \opt{Line} and coordinates of your own, with names and highlighted tensors:
\example{diagramcircle}
\example{diagramline}
\example{diagramcoordinates}
\keepfor[4]{diagramsymbols}Tensors passed as symbols are named after them, and the dimensions are written:
\example{diagramsymbols}
\keepfor[5]{diagramdimensions}The path, with the dimensions given but not written, and with the default dimension of every index:
\example{diagramdimensions}
\example{diagramdefault}
\texttt{"Steps" -> True} gives a list of pictures with the same plot range, to be shown by \texttt{ListAnimate}:
\example{diagramsteps}

\section{EinsumTree}
\keepfor[12]{tree}\fnentry{EinsumTree}{\fn{EinsumTree[spec, t1, t2, \dots]} or \fn{EinsumTree[spec, \{dims1, dims2, \dots\}]} draws the
contraction tree of the path: the leaves are the tensors, with their subscripts; each box is an intermediate, with its
indices, the indices summed at that step (or ``batch'' and ``outer'' when none is summed) and its cost; the box of the
result has an orange border. The title gives the specification, the cost of the path, the naive cost, the speed-up and the
largest intermediate. Without dimensions (\fn{EinsumTree[spec]}) the path is the one for \opt{DefaultDimension} and the
costs are not written. Options \opt{TensorNames}, \opt{Path}, \opt{DefaultDimension}, \texttt{ImageSize} and
\texttt{Background}, as in \fn{EinsumDiagram}.}
\example{tree}
\example{treenodims}

\section{The plan cache}
\fnentry{EinsumClearCache}{\fn{Einsum} keeps the plan of each call, keyed by the specification, the shapes and the options
\opt{Path} and \opt{BatchMethod}; a second call with the same key goes straight to the kernels. The cache holds up to 2000
plans and is emptied when full; \fn{EinsumClearCache[]} empties it.}
\example{clearcache}

\section{Messages}
\fn{Einsum}, \fn{EinsumPath}, \fn{EinsumFunction}, \fn{EinsumDiagram} and \fn{EinsumTree} give a message and
\texttt{\$Failed} for an invalid input; \texttt{Einsum::optbig} and the message of an unknown option are warnings, and the
computation goes on.

{\small\begin{tabularx}{\linewidth}{lX}\toprule message & when \\ \midrule
\texttt{Einsum::spec} & the specification cannot be read \\
\texttt{Einsum::nargs} & the number of tensors is not that of the subscripts \\
\texttt{Einsum::rank} & a tensor has not as many axes as its subscript has indices \\
\texttt{Einsum::ellrank} & a tensor has fewer axes than the indices named next to its ellipsis \\
\texttt{Einsum::ell} & a subscript has more than one ellipsis \\
\texttt{Einsum::ellout} & the operands have ellipsis axes, and the explicit output has no ellipsis \\
\texttt{Einsum::dim} & an index has different dimensions (other than 1) \\
\texttt{Einsum::outrep} & an index repeats in the output \\
\texttt{Einsum::outnf} & an index of the output is in no operand \\
\texttt{Einsum::path} & the path is not one of the strategies or a complete list of pairs \\
\texttt{Einsum::array} & an operand is a list that is not a rectangular array \\
\texttt{Einsum::mismatch} & an \fn{EinsumFunction} gets tensors of other dimensions (section~5) \\
\texttt{Einsum::optbig} & \opt{Optimal} is asked for more than 14 tensors: \opt{Greedy} is used \\
\texttt{Einsum::optx} & an unknown option (section~3) \\
\texttt{SummaUpdateCheck::latest} & the latest release cannot be read from GitHub \\ \bottomrule\end{tabularx}}

\smallskip
\example{errspec}
\example{errnargs}
\example{errrank}
\example{errellrank}
\example{errell}
\example{errellout}
\example{errdim}
\example{erroutrep}
\example{erroutnf}
\example{errpath}
\example{errarray}
\example{erroptbig}

\section{Performance}
\begin{itemize}
\item \textbf{Packed arrays.} \texttt{Developer\textasciigrave ToPackedArray}, \texttt{RandomReal} and \texttt{N} give packed arrays of
  machine numbers, on which every kernel is fast; \fn{Einsum} packs its operands when it can, and its results on
  machine numbers come out packed:
\end{itemize}
\example{packed}
\begin{itemize}
\item \textbf{Loops.} In a loop, use \fn{EinsumFunction}, or call \fn{Einsum} always with the same shapes, so that the plan
  comes from the cache.
\item \textbf{Large networks.} Beyond 8 tensors the default path is \opt{Greedy}. When the planning time is negligible
  against the contraction, try \texttt{"Path" -> "Optimal"} (up to 14 tensors), and compare the two with \fn{EinsumPath}.
\item \textbf{Batched contractions.} The cost model of \opt{BatchMethod} uses estimates, not timings of your machine; a kernel
  can be forced with \texttt{"BatchMethod" -> "Dot"} or \opt{Vectorized}.
\item \textbf{Sparse arrays.} A \texttt{SparseArray} stays sparse in the steps that are matrix products (\texttt{Dot}); the other
  kernels make it dense.
\item \textbf{Timings.} \texttt{Examples/Benchmarks.wl} measures \fn{Einsum} and \fn{EinsumFunction} on your machine against
  \texttt{Dot}, \texttt{MapThread}, \texttt{TensorContract} and \texttt{Total}.
\end{itemize}

\section{Limitations}
\begin{itemize}
\item An index cannot repeat in the output (\texttt{"i->ii"} is not allowed, as in numpy).
\item The exact search of the path is exponential in the number of tensors: the default is \opt{Greedy} beyond 8 tensors, and
  \opt{Optimal} is accepted up to 14.
\item A \texttt{SparseArray} stays sparse only in the steps that are matrix products.
\item The constants of the cost model of the batched contractions are estimates, not calibrated on your machine.
\end{itemize}

\section{Citing Summa}
If you use Summa in published work, please cite it as described in \texttt{CITATION.cff}:
\begin{code}
@software{Summa,
  author  = {Africani, Nazzareno},
  title   = {Summa: fast Einstein summation (einsum) for the Wolfram Language},
  version = {1.0.0},
  year    = {2026},
  url     = {https://github.com/n4sd4q/Summa}
}
\end{code}

\section{License}
Summa is free software, released under the GNU General Public License, version 3 or later
(\href{https://github.com/n4sd4q/Summa/blob/main/LICENSE}{\texttt{LICENSE}}), with two additional terms under its
section~7 (\href{https://github.com/n4sd4q/Summa/blob/main/NOTICE}{\texttt{NOTICE}}): the attribution ``Summa, by Nazzareno
Africani'' and the contents of \texttt{NOTICE} must be kept in all copies and modified versions, in their documentation and in
the message printed when the program is loaded in a notebook; and a modified version must be marked as different from the
original: it must not be distributed under the name ``Summa'' alone, and its documentation and its loading message must say
that it is a modified version of Summa by Nazzareno Africani, by whom and when it was modified. Copyright \copyright{} 2026
Nazzareno Africani (SISSA, Trieste).

\end{document}
